\chapter{结论与展望}

\section{研究结论}

本文从统计力学的视角系统研究了N皇后问题，将其映射为具有成对排斥相互作用的
格点气体模型，并通过大规模Monte Carlo模拟和张量网络精确枚举两条互补路径，
刻画了该系统的完整热力学行为。主要结论如下：

（1）建立了N皇后格点气体的统计力学框架。哈密顿量自然计数互相攻击的皇后对数，
基态（$E=0$）正好是$Q(N)$种非攻击解。每皇后基态熵$s_{0}=\ln N-\gamma$
（$\gamma\approx1.944$）可以逐级拆解：取每行放一个皇后为前提，单是禁止列重复
就减掉恰好1单位每皇后熵（源于Stirling公式$\ln N! = N\ln N - N$）；
再叠加两族对角线约束，又各自削去约$0.47$单位。
各约束的熵代价在大$N$下趋于常数，反映出系统在热力学极限下行为的高度规整。

（2）给出了高温极限的严格解析解。每皇后能量
$E/N|_{T\to\infty}=(5N-1)(N-1)/[3N(N+1)]$
对所有有限$N$精确成立，大$N$下趋于$5/3$；
高温熵$S(\infty)/N=\ln N+1+O(\ln N/N)$则给出了热力学积分的
高精度上界。Monte Carlo实测值在$T=500\,J$处与上述解析预言符合到万分之几，
交叉验证了推导与模拟两端。

（3）对$N=8\sim1024$八个尺寸、每温度点$10^{8}$次sweep的模拟数据显示：
每皇后比热$C_{v}/N$在$N\ge32$之后塌缩为同一条与$N$无关的曲线。
峰值$C_{v}^{\max}/N\approx1.63$位于$T^{*}\approx0.235\,J$，不随尺寸发散，
排除了热力学相变的可能。这个比热峰本质上是一个Schottky型异常——它来自
非攻击基态（$E=0$）与低能攻击构型之间有限能隙的热激发。借助哈密顿量的
攻击线分解，可严格证出最小能隙恰为一个耦合常数（$\Delta_{\min}=J$），
对应于在基态解上"翻转出一对攻击"的最低激发；实测峰位$T^{\ast}/J\approx 0.235$
比单纯两能级Schottky的理论峰位$0.42$更低，说明N皇后的激发态密度并非单一能级，
而是沿能量轴有丰富的多能级展宽。

（4）将$C_{v}/N$的收敛性与高温熵的精确值结合，通过热力学积分
$\int_{0}^{\infty}C_{v}/T\,dT$ 从Monte Carlo数据中独立提取了Simkin常数。
在$N=1024$处得到$\gamma_{\rm MC}=1.946\pm0.003$，与组合学结果
$\gamma=1.94400(1)$仅差$0.11\%$（约$0.002$），落在了jackknife误差
$\pm0.003$的范围之内。这构成了对$\gamma$的一条完全独立于组合枚举的物理检验。
$\gamma_{\rm MC}(N)$随$N$单调向$\gamma$靠近的趋势，则表明余下的偏差中
含有可辨认的有限尺寸修正。

（5）搭建了一套基于矩阵乘积算符的张量网络精确计数方案。四个方向的非攻击约束
被压缩进bond dimension $D=2$的MPO中，交汇后得到仅含17个非零元素的rank-9
局域张量。该张量网络对应的转移矩阵具有幂零性（$\mathcal{T}^{N+1}=0$），
这一点把它和常规统计力学模型从根本上区分开来——后者的配分函数由最大特征值
控制，而N皇后转移矩阵的所有特征值均为零，使得CTMRG、VUMPS等依赖谱分解的
标准无限尺寸方法在此失效；正确的做法是直接做有限$N$精确收缩。

\section{展望}

本工作也留下了一些待解决的问题。第一，如果把Monte Carlo模拟推到更大的$N$
（比如$2048$、$4096$乃至$8192$），热力学积分的精度会进一步提升，也更有可能
从数据中读出$Q(N)$渐近展开中有限$N$修正项的具体函数形式。有了它，目前的经验
标度外推就可以替换为有解析依据的预测性外推。

第二，方法论上的交叉值得尝试。计算量子物理和统计力学里最近发展起来的一些工具，
比如张量网络重正化群、神经网络变分Monte Carlo、用生成模型做重要性抽样等，
能不能用来处理N皇后这类经典组合问题并给出新的定量结果，目前还是一个开放
且相当有吸引力的问题。

第三，N皇后格点气体的热力学图景还有拓展空间。本文在正则系综下确认了$C_{v}/N$
的收敛性、排除了传统相变，但如果换到巨正则系综（允许皇后数目涨落），温度和
密度的二维参数空间里可能出现更复杂的相行为。一个自然的猜想是：不同几何方向
的约束（列、两族对角线）可能在不同温区依次"解冻"，在比热上留下有结构的
交叉序列，而不是现在看到的单个宽峰。

第四，张量网络收缩的指数瓶颈可以借助对称性部分缓解。N皇后问题的旋转和反射
对称性是现成的——把它们显式编入张量网络的结构中，可以削减状态空间的维数，
从而把可精确计算的$N$再往前推一步。

总的来说，本文的工作说明：N皇后问题远不止是一个经典组合学趣题。
作为一个约束几何规整的CSP范例，它天然坐落在统计力学、组合数学和计算物理
三者的交点上，可以充当交叉研究的实验田。用热力学手段测定Simkin常数是一条此前
未被尝试过的路，而这一思路的推广——利用有限温度模拟加积分来提取组合计数信息
——对其他计数问题或许同样管用。
